\SetPicSubDir{chap-appe}

\subsection{Office tracking and visualization technology adoption cases}
\label{app:Tracking}
\autoref{table:OccuTrack} shows examples of offices that utilize worker's location information during operation.

% \begin{table}[pos=hbt!]
%   \centering
%   \centering
\small
\setlength{\tabcolsep}{4pt}
\renewcommand{\arraystretch}{1.2}
\begin{tabularx}{\textwidth}{l p{3cm} p{4cm} p{6cm}}
\toprule
\textbf{Ref.} & \textbf{Company} & \textbf{Office} & \textbf{Technology} \\
\midrule
\cite{nihonsekkeiinc.RealtimeVisualization2024} & Nihon Sekkei & Head Office (Tokyo) & Smartphone + BLE beacons (Digital Twin visualization) \\
\cite{the_gear_kajima_2024} & Kajima Corporation & Singapore HQ ``The GEAR'' & IoT sensors + AI analytics\\
\cite{itokiITOKI2011} & Itoki & ITOKI TOKYO XORK & BLE beacons + mobile app \\
\cite{kokuyoco.OneYearRecords2021} & Kokuyo & Head Office & BLE beacons + mobile app \\
\cite{uchidayokoPressReleaseUchidaYoko2024} & Uchida Yoko & Shinkawa HQ Office, Olympus HQ, etc. & Wi-Fi indoor positioning + office map visualization \\
%\cite{colorkrew_biz_colorkrew_nodate} & Marubeni Corp. & HQ Chemical Division & QR code desk management (Colorkrew Biz) \\
%\cite{acall_it_2022} & Hashimoto Gumi & HQ Building & QR code desk management (Acall) \\
\cite{tanaka_development_2022} & Takenaka Corporation & Takenaka Central Bldg. South & AI cameras (human distribution \& posture sensing) \\
\bottomrule
\end{tabularx}

%   \caption{Office tracking and visualization technology adoption cases in Japan and Singapore}
%   \label{table:OccuTrack}
% \end{table}

\subsection{Occupancy Log Reconstruction Procedure}
\label{app:occupancy_reconstruction}

The raw data stream for occupancy tracking contains three event types from a face-recognition system:
\begin{itemize}[noitemsep]
    \item \textbf{Room Entry Events:} Triggered when a participant enters a specific room through a door-mounted access point.
    \item \textbf{General Presence Events:} Triggered when a participant is detected in shared zones or corridors, confirming their presence in the building.
    \item \textbf{Building Exit Events:} Triggered when a participant leaves the building through a main egress gate. Note that room-level exit events are not always recorded.
\end{itemize}

From these discrete events, a structured ``stay log'' is generated to represent continuous periods of presence. For each participant and each day, the following rules are applied:
\begin{itemize}[noitemsep]
    \item A stay begins at the timestamp of an “Entry” or equivalent presence-confirming event, and ends at the next recorded event (room entry, zone presence, or building exit).
    \item To prevent overnight carryover, any stay still active at the day’s final event is terminated at that time, typically corresponding to a building exit.
\end{itemize}

This process yields a stay log with anonymized IDs, locations, and start–end times for subsequent analysis.

\subsection{Occupancy tracking studies and examined parameters}
\label{app:OccupancyStudy}
Below \autoref{table:OccuMetrics} list up research studies that examined dynamic occupancy and the parameters used to measure dynamics.

\begin{table}[pos=htbp!]
  \centering
  \input{pic/chap-appe/ExOccuStudy}
  \caption{Summary of occupancy-related literature and metrics}
  \label{table:OccuMetrics}
\end{table}

\subsection{Thermal comfort survey contents}
\label{app:ThermalSurvey}
Below \autoref{table:comfsurvey} list up research studies that examined dynamic occupancy and the parameters used to measure dynamics.

\begin{table}[pos=htbp!]
  \centering
  \input{pic/chap-appe/ThermalComfort}
  \caption{Contents of Thermal comfort survey}
  \label{table:comfsurvey}
\end{table}


\SetPicSubDir{ch4-Results}
\subsection{Thermal comfort survey contents}
\begin{figure}[pos=htbp!]
    \centering
    \includegraphics[width=1.0\linewidth]{\Pic{png}{/agentic/Result_GEARsimulation/OperationExample}}
    \caption{Example of daily temperature control and predicted comfort performance both on \ac{sgcm} and \ac{rtgcm}}
    \label{fig:OpeExampls}
\end{figure}